\documentclass[10pt,a4paper]{amsart}
\usepackage[margin=19mm]{geometry}
\usepackage{amsmath,amssymb,mathtools}
\usepackage[scr=rsfs,cal=euler]{mathalfa}
\usepackage{xfrac}
\usepackage[unicode,hidelinks,bookmarks=false]{hyperref}
\newcommand{\ie}{i.e.\ }
\newcommand{\cf}{cf.\ }
\newcommand{\resp}{respectively\ }
\newcommand{\Subsection}{\subsection}
\providecommand{\nobreakdash}{\nobreak\hbox{-}}
\setlength{\emergencystretch}{2em}
\multlinegap0pt
\usepackage{bm}
\usepackage{bbm}
\usepackage{dsfont}
\usepackage{xspace}
\usepackage{cancel}
\usepackage{mathtools}
\usepackage{cleveref}
\usepackage{amsthm}
\makeatletter
\@ifundefined{theorem}{\newtheorem{theorem}{Theorem}}{}
\@ifundefined{lemma}{\newtheorem{lemma}[theorem]{Lemma}}{}
\makeatother
\title[Sticky CIR process with potential: invariant measure and exa]{Sticky CIR process with potential: invariant measure and exact sampling}
\author{Tony Shardlow}
\date{Source: arXiv:2605.13648v5 (CC BY 4.0)}
\begin{document}
\maketitle

\noindent\emph{Source and licence.} Sticky CIR process with potential: invariant measure and exact sampling. Version arXiv:2605.13648v5 (CC BY 4.0), distributed under the Creative Commons Attribution 4.0 International licence, as recorded in the arXiv licence metadata for this exact version and shown on the abstract page https://arxiv.org/abs/2605.13648v5. Full rights notice: licenses/arxiv-2605.13648-CC-BY-4.0.txt.\par\smallskip

\noindent\textit{Editorial scope.} Counted unit: the Theorem whose source label is \texttt{thm:transition-dist} in the article (source0.tex lines 837--868, source label \texttt{thm:transition-dist}), delivered together with the article's own complete original proof of it (source0.tex lines 870--926, one \texttt{proof} environment of 57 lines). No mathematics is written, repaired or re-derived: statement and proof are the article's own text line for line. Required context retained uncounted: eq:green-resolvent (lines 612--614); eq:wentzell (lines 732--742); lem:sticky-resolvent (lines 798--816); eq:cmu (lines 806--808). Every definition, display and label of those lines is retained unchanged. Omitted from this package and not redistributed: 1--611, 615--731, 743--797, 809--836, 869--869, 927--2464. Nothing inside a retained excerpt is omitted or altered: every retained body is delivered exactly as the article's lines are written, so no omission receipt of any kind accompanies this package. Citations: the retained excerpts carry no citation command at all, so no bibliography entry of the article is transcribed and no reference is redistributed. Reference adaptation: none was needed and none was made; every reference inside the delivered unit resolves inside it, so no unresolved reference or citation is produced. Printed slips kept literal: no printed word, symbol, hypothesis or number of the counted statement or of its proof was altered. Layout and class: the article's own class is \texttt{imsart} with packages including fontenc, amsthm, amsmath, natbib, hyperref, mathtools, amsfonts, ifdraft, stmaryrd, microtype, amssymb, booktabs, enumitem, inputenc; the wrapper is the lane's pinned \texttt{amsart} editorial wrapper at 10pt on A4, which restates the theorem-like environments the retained text needs and adds the article's own macro definitions that the retained text needs (none), with extra wrapper packages (bm, bbm, dsfont, xspace, cancel, mathtools, cleveref). The delivered numbering is the unit's own. Assets: no retained span contains a figure environment, a graphics-inclusion command, a PDF-page inclusion, a table or a TikZ picture; no image file is copied into this package, the package declares no asset dependency and no image file is redistributed with it. Licence: version arXiv:2605.13648v5 of this article is distributed under the Creative Commons Attribution 4.0 International licence, as recorded in the arXiv licence metadata for this exact version and shown on the abstract page https://arxiv.org/abs/2605.13648v5. A full rights notice is delivered at licenses/arxiv-2605.13648-CC-BY-4.0.txt. Characters: every retained excerpt is pure ASCII, so the delivered engine, XeLaTeX with its default Latin Modern text font, reports no missing character.
\begin{equation}\label{eq:green-resolvent}
  G_\alpha(x,y)\,m(dy) = \frac{1}{\alpha}\,\mathbb{P}(\tilde u_{T_\alpha}\in dy\mid \tilde u_0=x),
\end{equation}

\begin{lemma}[resolvent boundary condition]\label[lemma]{lem:resolvent-bc}
The resolvent kernel $G^{\rm sticky}_\alpha(x,y)$, viewed as a function
of its first argument at fixed $y \in (0,\infty)$, solves
$(\alpha-\mathcal{L})G^{\rm sticky}_\alpha(\cdot,y)=\delta_y$ on $(0,\infty)$
in the speed-measure sense, with boundary condition
\begin{equation}\label{eq:wentzell}
  \frac{dG^{\rm sticky}_\alpha(\cdot,y)}{ds}(0^+)
  = \frac{\alpha}{\mu}\,G^{\rm sticky}_\alpha(0,y).
\end{equation}
By the symmetry $G^{\rm sticky}_\alpha(x,y) = G^{\rm sticky}_\alpha(y,x)$, the same condition holds in the second argument.
\end{lemma}

\begin{lemma}[sticky resolvent kernel]\label[lemma]{lem:sticky-resolvent}
The fundamental solutions of $(\alpha-\mathcal{L})f=0$ satisfying the
boundary condition~\eqref{eq:wentzell} are $f_\infty(x) = U(a,b,z_x)$
and
\begin{equation}\label{eq:f0sticky}
  f_0^{\rm sticky}(x) = M(a,b,z_x) + c_\mu\,U(a,b,z_x),
\end{equation}
where
\begin{equation}\label{eq:cmu}
  c_\mu = \frac{-\alpha}{\mu\,|\mathcal{W}| + \alpha U_0}.
\end{equation}
The sticky resolvent kernel is
\[
  G^{\rm sticky}_\alpha(x,y) =
  \frac{f_0^{\rm sticky}(x\wedge y)\,f_\infty(x\vee y)}{|\mathcal{W}|},
\]
with the same normalisation constant $|\mathcal{W}|$ as the reflecting
resolvent kernel.
\end{lemma}

\begin{theorem}[transition distribution]\label{thm:transition-dist}
  For $x>0$, the transition distribution $P_x(u_{T_\alpha}\in\cdot)$
  admits the resolvent representation
  \begin{equation}\label{eq:trans}
    \mathbb{P}(u_{T_\alpha}\in A\mid u_0=x)
    = \frac{\alpha}{\mu}\,G^{\rm sticky}_\alpha(x,0)\,\mathbf{1}_{0\in A}
    + \alpha\int_{A\cap(0,\infty)} G^{\rm sticky}_\alpha(x,y)\,m'(y)\,dy,
  \end{equation}
  equivalent to the three-component mixture
  \begin{equation}\label{eq:trans-mixture}
    \mathbb P(u_{T_\alpha}\in A\mid u_0=x)
    = w_0(x)\,\mathbf{1}_{0\in A}
    + w_<(x)\,\nu_<(A)
    + w_>(x)\,\nu_>(A),
  \end{equation}
  where $\nu_<$ and $\nu_>$ are the probability measures with Lebesgue
  densities
  \[
    \frac{d\nu_<(y)}{dy}  \propto  f_0^{\rm sticky}(y)\,m'(y)\,\mathbf{1}_{(0,x)}(y),
    \qquad
    \frac{d\nu_>(y)}{dy}  \propto  U(a,b,z_y)\,m'(y)\,\mathbf{1}_{(x,\infty)}(y),
  \]
  and the weights are
  \begin{align}
    w_0(x) &= -c_\mu\,U(a,b,z_x), \label{eq:w0}\\
    w_<(x) &= \frac{\alpha\,U(a,b,z_x)}{|\mathcal{W}|}
               \int_0^x f_0^{\rm sticky}(y)\,m'(y)\,dy, \label{eq:wlt}\\
    w_>(x) &= \frac{\alpha\,f_0^{\rm sticky}(x)}{|\mathcal{W}|}
               \int_x^\infty U(a,b,z_y)\,m'(y)\,dy, \label{eq:wgt}
  \end{align}
  satisfying $w_0(x)+w_<(x)+w_>(x)=1$.
\end{theorem}

\begin{proof}
  The representation~\cref{eq:trans} follows from~\cref{eq:green-resolvent}
  and the decomposition of the speed measure,
  $m(dy)=\mu^{-1}\delta_0(dy)+m'(y)\,dy$.

  For~\cref{eq:trans-mixture}, evaluate~\cref{eq:trans} on the atom first.
  By~\cref{lem:sticky-resolvent} with $y=0$,
  $G^{\rm sticky}_\alpha(x,0) = f_0^{\rm sticky}(0)\,U(a,b,z_x)/|\mathcal{W}|$.
  Evaluating $f_0^{\rm sticky}(0) = M(a,b,0)+c_\mu U(a,b,0) = 1+c_\mu U_0$
  and substituting $c_\mu=-\alpha/(\mu|\mathcal{W}|+\alpha U_0)$
  from~\cref{eq:cmu} gives
  \begin{equation}\label{eq:f0sticky-zero}   
    f_0^{\rm sticky}(0)
    = 1 - \frac{\alpha U_0}{\mu|\mathcal{W}|+\alpha U_0}
    = \frac{\mu|\mathcal{W}|}{\mu|\mathcal{W}|+\alpha U_0}.
  \end{equation}

  The atom weight $w_0(x) = (\alpha/\mu)\,G^{\rm sticky}_\alpha(x,0)$ then
  becomes
  \[
    w_0(x)
    = \frac{\alpha}{\mu}
      \,\frac{\mu|\mathcal{W}|}{\mu|\mathcal{W}|+\alpha U_0}
      \,\frac{U(a,b,z_x)}{|\mathcal{W}|}
    = -c_\mu\,U(a,b,z_x) > 0,
  \]
  where the inequality holds since $c_\mu<0$ and $U(a,b,z_x)>0$.

  To verify $w_0+w_<+w_>=1$, recall that
  \[
    G_\alpha(x,y) = \frac{M(a,b,z_{x\wedge y})\,U(a,b,z_{x\vee y})}{|\mathcal{W}|}
  \]
  is the reflecting resolvent kernel (the case $c_\mu=0$), and use the decomposition
  \[
    G^{\rm sticky}_\alpha(x,y)
    = G_\alpha(x,y)
    + \frac{c_\mu\,U(a,b,z_x)\,U(a,b,z_y)}{|\mathcal{W}|},
  \]
  which follows from $f^{\rm sticky}_0 = M + c_\mu U$. Since
  $w_<+w_> = \alpha\int_0^\infty G^{\rm sticky}_\alpha(x,y)\,m'(y)\,dy$,
  the decomposition gives
  \[
    w_<+w_>
    = \alpha\int_0^\infty G_\alpha(x,y)\,m'(y)\,dy
    + \frac{\alpha c_\mu\,U(a,b,z_x)}{|\mathcal{W}|}\int_0^\infty U(a,b,z_y)\,m'(y)\,dy.
  \]
  The first term equals $1$ since $\alpha G_\alpha(x,\cdot)\,m'(\cdot)$
  is a probability density on $(0,\infty)$ (the boundary mass vanishes for
  the reflecting process; see \cref{eq:green-resolvent}).
  The second integral equals $|\mathcal{W}|/\alpha$ by the same argument at
  $x=0$, where $G_\alpha(0,y)=U(a,b,z_y)/|\mathcal{W}|$; that is,
  \begin{equation}\label{eq:U-integral}
    \int_0^\infty U(a,b,z_y)\,m'(y)\,dy = |\mathcal{W}|/\alpha,
  \end{equation}
  so the second term equals $c_\mu\,U(a,b,z_x)=-w_0$.
  Hence, $w_0+w_<+w_>=1$.
\end{proof} 

\end{document}
